\project{01}{LTE Cat-4 motion-triggered gateway}{Raspberry Pi 4 + SIM7600E-H + ADXL345}{the only Cat-4 default route, plus shock-event compression}

\begin{unique}
Owns the only Cat-4 \emph{default route} in the book, plus ADXL345 event compression. Cat-M and NB-IoT are P03 and P02.
\end{unique}

\begin{blueprint}
Draft P5 is the same radio used as a NAT router; draft P1 used Intel MKL. We use QMI/RNDIS for the WAN and NumPy for nothing heavier than a magnitude check. Do not also put GNSS tracking here, that is P03.
\end{blueprint}

\subsection*{Intent}

The Pi 4 gets a default IPv4 route through the SIM7600E-H over USB, as RNDIS or as QMI-WWAN. The ADXL345 publishes a shock event only when the acceleration magnitude leaves a dead-band around 1 g, so the uplink carries events rather than a sample stream. Three LEDs report the three states that matter on a bench: the modem registered, the default route moved to the modem, and an event fired.

The lab is a gateway, not a tracker and not a router. Two things are deliberately out of scope. GNSS belongs to P03, because a Cat-4 module that also chases a fix hides which subsystem failed. Network address translation for other hosts belongs to nothing in this book, because one SIM and one bench do not need it.

\begin{kit}
Pi 4, SIM7600E-H with LTE and GNSS antennas and a micro-SIM, SEN0032 ADXL345, three LK-LED10 modules on BCM 16/20/21, a 3 A USB-C supply, keyboard.
\end{kit}

\subsection*{System architecture}

\diagram{p01_arch}{The two planes of the lab. The user plane is the modem's USB network interface; the control plane is an AT character device on the same cable. The accelerometer never touches either.}

The SIM7600E-H presents several interfaces on one USB cable. One of them is a network device that carries the user plane, and another is a character device that accepts AT commands. Keeping those apart is the whole discipline of the lab: the WAN comes up through \texttt{dhcpcd} or ModemManager on the network device, and diagnostics go to the AT device. Point-to-point protocol is not used at all, because running it alongside QMI or RNDIS gives two things one modem cannot serve at once.

The event path is independent of the network path. A Python loop reads six registers over I2C at 20 Hz, computes a magnitude, and calls \texttt{mosquitto\_pub} when the magnitude leaves the dead-band. If the WAN is down the publish fails and the loop continues, which is the behaviour you want from a field gateway.

\subsection*{Wiring and schematic}

\diagram{p01_schematic}{The accelerometer on the HAT pass-through header and the three status LEDs. SDO left floating selects address 0x53.}

\begin{table}[H]\centering\small
\begin{tabular}{p{42mm}p{26mm}p{22mm}p{58mm}}
\toprule
Signal & Pi header pin & BCM GPIO & Note \\
\midrule
ADXL345 VCC & 1 & 3V3 & Never 5 V \\
ADXL345 SDA & 3 & GPIO2 & I2C1 data \\
ADXL345 SCL & 5 & GPIO3 & I2C1 clock \\
ADXL345 GND & 6 & GND & \\
ADXL345 SDO & & & Left floating, selects 0x53 \\
ADXL345 INT1 & 11 & GPIO17 & Optional, not needed for the magnitude poll \\
Green module, S1 & 36 & GPIO16 & CEREG registered. Lit when driven high \\
Yellow module, S1 & 38 & GPIO20 & Default route on usb0 or wwan0 \\
Red module, S1 & 40 & GPIO21 & Last shock event \\
Module grounds & 34 or 39 & GND & The breadboard ground rail, as in P15 \\
Pins 13, 15, 16 & & GPIO27, 22, 23 & Not free on this HAT. See the note below \\
\bottomrule
\end{tabular}
\caption{Wiring. The HAT carries a full 40-pin pass-through, so the accelerometer and the modem coexist on one header and pins 3 and 5 are untouched by the HAT. The indicator pins are not the ones the draft used; the note below says why.}
\end{table}

\begin{note}{The status triple cannot sit on BCM 27, 22 and 23 under this HAT}
The draft put the three indicators on header pins 13, 15 and 16. On the SIM7600E-H those are the modem's own control lines: GPIO27 is the ring indicator, GPIO22 is data terminal ready and GPIO23 is clear to send. An LED hung on any of them is a second driver on a line the modem is already using, and the failure it produces is not a dark LED but a modem that behaves oddly for reasons nothing in the log explains.

This is not peculiar to one board. On the SIM7020E of P02, GPIO27 is the module's power key, which the vendor scripts pulse to boot it. An indicator there would switch the modem on and off every time it reported a change of state, so the lab would interfere with the thing it is reporting.

So the triple is per lab rather than per book. The pins above are chosen from the high end of the header, where neither modem HAT documents a connection, with the module grounds going to pin 34 or 39. \textbf{Confirm them against the pinout of the HAT in front of you before wiring}, because this is a fact about one board and not about the Raspberry Pi. Some Waveshare HATs carry a jumper block that selects which Pi GPIO reaches each control line, and on a board that has one, removing three jumpers is the other way to free the original pins.
\end{note}

\begin{asciiart}
ADXL345                 Pi 40-pin (HAT pass-through pins 1/3/5/6 still available)
VCC ------------------- pin 1  3V3      NEVER 5V
GND ------------------- pin 6  GND
SDA ------------------- pin 3  GPIO2
SCL ------------------- pin 5  GPIO3
SDO floating => 0x53
INT1 optional --------- BCM17  (not required for the magnitude poll)

SIM7600 jumpers:
  VCCIO = 3V3
  PWR   = 3V3   (auto-on). Use PWR=D6 only if you want a GPIO power-key.
  UART  = B if you need ttyAMA0; prefer the HAT USB cable for AT + data.
USB cable: HAT USB-A/micro -> Pi USB-A  (user plane + AT channels)
Antennas: LTE main + GNSS. No antenna = no attach, and can stress the PA.
\end{asciiart}

\subsection*{Bench layout}

\diagram{p01_bench}{Bench layout. The short USB cable between the HAT and the Pi carries the user plane, so it is not optional. The 3 A supply is not optional either: a transmit burst on Cat-4 draws current the Pi cannot lend.}

Seat the HAT, fit both antennas and insert the SIM before applying power. Connect the HAT USB cable in the same pass. A modem that is powered with no antenna fitted can stress its own power amplifier, and a SIM inserted under power is not detected until the next boot.

\subsection*{Software design (UML)}

\diagram{p01_uml}{Bring-up sequence. The AT channel confirms registration, the network interface carries traffic, and the accelerometer loop is a separate process that survives a WAN outage.}

The publisher is a single loop with one piece of state, a cool-off timestamp. It is a token bucket in its simplest form: at most one event every two seconds, no matter how long the table keeps ringing. Without it, one fist produces a burst of twenty publishes and the value of the compression is lost.

\subsection*{Data flow (ASCII)}

\begin{asciiart}
  Pi 4                                            SIM7600E-H HAT
  +---------------------------------+             +------------------------+
  | p01_shock.py                    |             |  LTE Cat-4 radio       |
  |   i2c-1 -> 0x53 ADXL345 @20 Hz  |             |    |                   |
  |   |a| - 1g > 0.35 ?             |             |    v                   |
  |   |  yes -> mosquitto_pub ------------------> | usb0 / wwan0  (user)   |---> broker
  |   |  cool-off 2 s               |             |                        |
  |   v                             |   USB       | /dev/ttyUSB2  (AT) <--------+
  | gpio: GRN 16  YEL 20  RED 21    |<===========>|                        |     |
  +---------------------------------+             +------------------------+     |
        ^                                                                        |
        +--- AT+CSQ every 60 s, AT+CEREG? for the green LED --------------------+
\end{asciiart}

\subsection*{Steps}

\begin{steps}
\item \textbf{Prepare the host.} Flash Raspberry Pi OS Bookworm 64-bit. Enable I2C under \texttt{raspi-config}, Interface options. Use a 3 A supply.

\item \textbf{Assemble before power.} Seat the HAT, fit both antennas, insert the SIM, and connect the HAT USB cable, all before applying power.

\item \textbf{Identify the modem and the accelerometer.}
\begin{shellcode}
lsusb | grep -iE "sim|1e0e"
dmesg | tail -n 40
ls /dev/ttyUSB* /dev/ttyACM* 2>/dev/null
sudo apt-get update
sudo apt-get install -y minicom i2c-tools mosquitto-clients python3-smbus python3-rpi.gpio
sudo i2cdetect -y 1          # must show 53
\end{shellcode}

\item \textbf{Open the AT port.} It is commonly \texttt{/dev/ttyUSB2} on the SIM7600. If minicom shows nothing, try USB3 then USB1.
\begin{atcode}
sudo minicom -D /dev/ttyUSB2 -b 115200
AT
ATE1
AT+CPIN?
AT+CSQ
AT+COPS?
AT+CEREG?
AT+CGDCONT?
\end{atcode}

\item \textbf{Bring up the WAN.} Prefer RNDIS or QMI over the point-to-point protocol.
\begin{shellcode}
# RNDIS path: usb0 typically appears when the HAT USB cable is connected
ip link
sudo dhcpcd usb0 || sudo dhclient usb0
ip route
# If usb0 did not appear, load qmi_wwan / option and use ModemManager:
sudo apt-get install -y modemmanager libqmi-utils
mmcli -L
# Do not run PPP in parallel with QMI or RNDIS.
\end{shellcode}

\item \textbf{Confirm and lock the default route}, so the modem wins over wlan0.
\begin{shellcode}
ip route | grep default
# optional: give wlan0 a higher metric so LTE wins
sudo ip route del default dev wlan0
\end{shellcode}

\item \textbf{Run the event publisher}, saved as \texttt{/home/pi/p01\_shock.py}.
\begin{pycode}
import time, math, json, subprocess, smbus
bus = smbus.SMBus(1); A = 0x53
bus.write_byte_data(A, 0x2D, 0x08)   # measure
bus.write_byte_data(A, 0x31, 0x0B)   # full-res +/-16g

def g(lo, hi):
    v = (hi << 8) | lo
    if v & 0x8000:
        v -= 65536
    return v / 256.0

cool = 0
while True:
    b = bus.read_i2c_block_data(A, 0x32, 6)
    ax, ay, az = g(b[0], b[1]), g(b[2], b[3]), g(b[4], b[5])
    mag = math.sqrt(ax*ax + ay*ay + az*az)
    if abs(mag - 1.0) > 0.35 and time.time() > cool:
        payload = json.dumps({"mag": round(mag, 3), "ax": round(ax, 3)})
        subprocess.run(["mosquitto_pub", "-h", "test.mosquitto.org",
                        "-t", "lab/p01/shock", "-m", payload], check=False)
        cool = time.time() + 2
    time.sleep(0.05)
\end{pycode}

\item \textbf{Drive the LEDs from state, not from hope.} Green follows \texttt{AT+CEREG?}, yellow follows the presence of a default route on \texttt{usb0}, red follows the last publish. Poll \texttt{AT+CSQ} every 60 s and log it.
\end{steps}

\subsection*{Acceptance test}

\begin{itemize}
\item A fist on the table produces exactly one MQTT JSON message and lights the red LED.
\item \texttt{ping -I usb0 8.8.8.8} succeeds, or the operator's own DNS address does.
\item \texttt{ip route | grep default} names the modem interface, not wlan0.
\item Outdoors, \texttt{AT+CGPS=1} then \texttt{AT+CGPSINFO} returns non-blank fields. That is a side check on the hardware, not this lab's product.
\end{itemize}

\subsection*{Practices}

User plane on USB, AT on a ttyUSB. Token-bucket the publisher, which is what the two-second cool-off is. Log \texttt{AT+CSQ} every 60 s so a weak-signal night is visible the next morning. Do not seat the MCC 118 or the 3.5 inch LCD at the same time: one 40-pin HAT per host.

\subsection*{Pitfalls}

\begin{itemize}
\item \textbf{No antenna fitted.} The module will not attach, and transmitting into an open port stresses the power amplifier. Fit both antennas before power.
\item \textbf{A 2.5 A supply.} The Pi 4 alone is close to that budget. A Cat-4 transmit burst on top of it produces undervoltage warnings in \texttt{dmesg} and random resets that look like software faults.
\item \textbf{Running the point-to-point protocol next to QMI.} Two user planes, one radio. The symptom is a route that appears and then carries nothing.
\item \textbf{Picking the wrong AT device.} Several \texttt{ttyUSB} nodes appear. The wrong one is silent rather than an error, so try USB2, then USB3, then USB1.
\item \textbf{Publishing every sample.} Without the cool-off, one impact produces a burst and the dead-band logic buys nothing.
\item \textbf{5 V on the accelerometer.} The SEN0032 goes to pin 1, never pin 2.
\end{itemize}

\subsection*{Sources}

\begin{itemize}
\item Waveshare SIM7600E-H 4G HAT wiki, \url{https://www.waveshare.com/wiki/SIM7600E-H_4G_HAT}
\item SIMCom SIM7600 series AT command manual (module vendor documentation)
\item DFRobot SEN0032 ADXL345 product wiki, \url{https://wiki.dfrobot.com/}
\item Analog Devices ADXL345 data sheet, register map 0x2D, 0x31, 0x32
\item ModemManager and libqmi documentation, \url{https://www.freedesktop.org/wiki/Software/ModemManager/}
\item Raspberry Pi hardware documentation, 40-pin header, \url{https://www.raspberrypi.com/documentation/computers/raspberry-pi.html}
\end{itemize}
